%=======================================================================
% 中文版标题页 — RetNet-0.5B
% 标题必须与 docs/26_RetNet-0.5B_技术报告.md 的 H1 一致（H1 不进正文）
% 摘要与目录由 mktex.py 紧随其后插入（摘要取自 Markdown 的 `## 摘要` 一节）
%=======================================================================
\thispagestyle{empty}

\begin{center}
\vspace*{2.6cm}

{\LARGE\bfseries RetNet-0.5B：逐段单卡训练到线上的\\[0.25em] 5.04 亿参数保留网络}

\vspace{0.45cm}

{\large\bfseries ——宽度生长为何未能解锁}

\vspace{0.7cm}

{\small RetNet-0.5B: Training a 504M-Parameter Retentive Network to Production,\\
and Why Width Growth Never Unlocked}

\vspace{0.8cm}

{\normalsize RetNet 项目组}

\vspace{0.25cm}

{\small\itshape 2026-10-13}

\vspace{1.6cm}

{% 表 1 的题注在正文里属于封面块（不进正文），因此手写在这里，编号与正文衔接
\small\textbf{表 1.} 模型摘要（数值出处见正文对应章节）}

\vspace{0.5em}

\begin{minipage}{0.94\textwidth}
\centering\small
\begin{tabular}{@{}rl@{}}
\toprule
模型 & 504.18\,M 参数（504,179,896），20 层，1408 维，22 头 \\
训练 & 三段各一个 epoch，实际消耗 1,656,209,408 token（1.656\,B） \\
硬件 & RTX 3090 24\,GB / RTX PRO 6000 96\,GB / RTX 3070 8\,GB \\
数据 & \code{chat\_pack\_v13}，552,330,489 token 槽位，2048 token 上下文 \\
结果 & 验证 loss 2.6358 $\rightarrow$ 1.9459；生长新增的 82 个仿射张量中 78 个逐位为 0 \\
\bottomrule
\end{tabular}
\end{minipage}

\vfill

{\small\itshape 本 PDF 为中文版；英文版为另一文件 RetNet-0.5B-Technical-Report.pdf。}

\vspace{1.0cm}
\end{center}

\clearpage
